\chapter*{Introducere}
\addcontentsline{toc}{chapter}{Introducere}

\section*{Motivație}
\addcontentsline{toc}{section}{Motivație}

În contextul actual al tehnologiei audio, flexibilitatea și calitatea procesării semnalului au devenit cerințe fundamentale atât în mediul profesional, cât și în cel de larg consum. Trecerea de la procesarea pur analogică la cea digitală a deschis noi oportunități pentru optimizarea răspunsului în frecvență, corecția acustică a spațiilor și implementarea unor algoritmi complexi de filtrare și dinamică.

Motivația acestui proiect rezidă în necesitatea de a crea un sistem audio versatil care să combine performanța etajelor analogice de precizie cu puterea de calcul a unui procesor digital de semnal (DSP). Un astfel de sistem permite utilizatorului să adapteze comportamentul amplificatorului în timp real, oferind o soluție compactă și eficientă pentru aplicații multimedia avansate.

\section*{Obiective}
\addcontentsline{toc}{section}{Obiective}

Obiectivul principal al lucrării este proiectarea, simularea și implementarea practică a unui amplificator audio multicanal dotat cu un procesor digital de semnal programabil. Pentru atingerea acestui scop, au fost stabilite următoarele obiective secundare:
\begin{itemize}
    \item Proiectarea unor interfețe de intrare compatibile cu standardele profesionale (+4dBu) și consumer (-10dBV), utilizând circuite balansate și nebalansate.
    \item Implementarea unui etaj de preamplificare cu zgomot redus bazat pe amplificatoare operaționale SoundPlus\texttrademark.
    \item Integrarea procesorului de semnal ADAU1701 și configurarea acestuia pentru procesare în timp real la 48kHz.
    \item Realizarea unui sistem de alimentare eficient, capabil să susțină atât etajele analogice de mic semnal, cât și etajele de putere în Clasa D.
    \item Dezvoltarea unui programator USB pentru configurarea facilă a algoritmilor DSP prin interfața SigmaStudio.
    \item Validarea performanțelor sistemului prin simulări software și măsurători experimentale.
\end{itemize}

\section*{Aplicabilitate}
\addcontentsline{toc}{section}{Aplicabilitate}

Sistemul propus își găsește aplicabilitatea în diverse domenii ale ingineriei audio și producției multimedia:
\begin{enumerate}
    \item \textbf{Sisteme de monitorizare audio:} Permite calibrarea precisă a boxelor de monitorizare în funcție de acustica camerei.
    \item \textbf{Prototipare rapidă:} Utilizarea DSP-ului programabil facilitează testarea rapidă a unor noi algoritmi de crossover, egalizare sau compresie fără modificarea hardware-ului.
    \item \textbf{Instalații audio personalizate:} Flexibilitatea intrărilor și ieșirilor face sistemul ideal pentru instalații fixe în spații comerciale sau rezidențiale unde este necesară o gestionare multicanal a semnalului.
\end{enumerate}
